\subsection*{5.2. The hard Lefschetz theorem (LV)}

\subsubsection*{5.2.1.}
Let $S$ be projective of dimension $n$ over the algebraically closed field
$k$, and let $S\hookrightarrow\mathbb P$ be a projective embedding.  Denote by
\[
u\in H^2(S,\mathbb Z_\ell(1))
\]
the cohomology class of a hyperplane section.  Thus $u$ is the pullback along
$S\hookrightarrow\mathbb P$ of the first Chern class of
$\underline{\mathcal O}_{\mathbb P}(1)$ in
$H^2(\mathbb P,\mathbb Z_\ell(1))$.  Write
\[
L_S:H^\bullet(S,\mathbb Q_\ell(j))
   \longrightarrow H^{\bullet+2}(S,\mathbb Q_\ell(j+1))
\]
for multiplication by $u$.

\subsubsection*{5.2.2.}
\underline{We say that $S$ (equipped with $u$) satisfies (LV), the hard
Lefschetz theorem, if, for every (or for one) integer $j$ and every integer
$i>0$, the maps}
\[
L_S^i:H^{n-i}(S,\mathbb Q_\ell(j))
\longrightarrow H^{n+i}(S,\mathbb Q_\ell(i+j))
\]
\underline{are isomorphisms.}

\subsubsection*{5.2.2.1.}
It is conjectured that this condition holds when $S$ is smooth.  It is true
when $k$ has characteristic zero [6].  The specialization theorem
[SGA 4, XVI 2.2] therefore gives the same conclusion when $S$, together with
its projective embedding, lifts to characteristic zero.  In particular, it
holds for a smooth complete intersection.  It is also known for a smooth $S$
of dimension at most two [2], and for an abelian variety, either directly
(Lieberman [4] and [2]) or by Mumford's theorem [5] that a polarized abelian
variety lifts, together with its polarization, to characteristic zero.

Recall that, for a projective variety $S$ and $i\leq n$, the primitive part of
cohomology is defined by
\begin{equation}\tag{5.2.2.2}
\operatorname{Prim}^i(S,\mathbb Q_\ell(j))
=\operatorname{Ker}\!\left(
L^{n-i+1}:H^i(S,\mathbb Q_\ell(j))
\longrightarrow
H^{2n-i+2}(S,\mathbb Q_\ell(j+n-i+1))
\right).
\end{equation}
If $S$ satisfies (LV), then for $i\leq n$ one has the decomposition
\begin{equation}\tag{5.2.2.3}
H^i(S,\mathbb Q_\ell(j))
\xleftarrow[\sim]{\ 1+L\ }
\operatorname{Prim}^i(S,\mathbb Q_\ell(j))
\oplus H^{i-2}(S,\mathbb Q_\ell(j-1)).
\end{equation}

\subsection*{Proposition 5.2.3.}
\underline{Let $X$ be smooth and projective and suppose that it satisfies
(LV).  Let $\tau:Y\hookrightarrow X$ be the inclusion of a smooth hyperplane
section.  Then}
\begin{equation}\tag{5.2.3.1}
\operatorname{Prim}^n(X,\mathbb Q_\ell(j))
=\operatorname{Ker}\!\left(
\tau^*:H^n(X,\mathbb Q_\ell(j))
\longrightarrow H^n(Y,\mathbb Q_\ell(j))
\right).
\end{equation}

\underline{Proof.}
The projection formula for $\tau:Y\to X$ gives the commutative diagram
\begin{equation}\tag{5.2.3.2}
\begin{tikzcd}[column sep=huge,row sep=huge]
H^n(X,\mathbb Q_\ell(j))
  \arrow[rr,"L"]\arrow[dr,"\tau^*"']
  & {} &
H^{n+2}(X,\mathbb Q_\ell(j+1))
  \arrow[dl,"\tau_*"]\\
{} & H^n(Y,\mathbb Q_\ell(j)) & {}
\end{tikzcd}
\end{equation}
By duality and weak Lefschetz,
\[
\tau_*:H^n(Y,\mathbb Q_\ell(j))
\longrightarrow H^{n+2}(X,\mathbb Q_\ell(j+1))
\]
is the transpose of an isomorphism and hence is itself an isomorphism.  The
assertion follows.

\subsection*{Lemma 5.2.4.}
\underline{Let $X$ be smooth and projective, and let
$\tau:Y\hookrightarrow X$ be the inclusion of a smooth hyperplane section.
Then $X$ satisfies (LV) if and only if the following two conditions hold:}

\smallskip
\noindent
\textup{5.2.4.1.} \underline{$Y$ satisfies (LV).}

\smallskip
\noindent
\textup{5.2.4.2.} \underline{The restriction of the cup product on
$H^{n-1}(Y,\mathbb Q_\ell)$ to the subspace
$\operatorname{Im}H^{n-1}(X,\mathbb Q_\ell)$ is nondegenerate.}

\underline{Proof.}
We have
\begin{equation}\tag{5.2.4.3}
\left\{
\begin{aligned}
L_X&=\tau_*\tau^*,\\[1ex]
L_Y&=\tau^*\tau_* \qquad\text{(by 4.2.11),}
\end{aligned}
\right.
\end{equation}
and consequently
\begin{equation}\tag{5.2.4.4}
L_X^{i+1}=\tau_*L_Y^i\tau^*.
\end{equation}
In the commutative diagram
\begin{equation}\tag{5.2.4.5}
\begin{tikzcd}[column sep=huge,row sep=huge]
H^{n-i-1}(X,\mathbb Q_\ell(j))
  \arrow[r,"L_X^{i+1}"]\arrow[d,"\tau^*"]
  &H^{n+i+1}(X,\mathbb Q_\ell(j+i+1))\\
H^{n-1-i}(Y,\mathbb Q_\ell(j))
  \arrow[r,"L_Y^i"]
  &H^{n-1+i}(Y,\mathbb Q_\ell(j+i))\arrow[u,"\tau_*"']
\end{tikzcd}
\end{equation}
$\tau^*$ is an isomorphism for $i\geq1$ by weak Lefschetz, and $\tau_*$ is
an isomorphism by duality.  Hence $X$ satisfies (LV) if and only if $Y$
satisfies (LV) and
\begin{equation}\tag{5.2.4.6}
L_X:H^{n-1}(X,\mathbb Q_\ell(j))
\longrightarrow H^{n+1}(X,\mathbb Q_\ell(j+1))
\end{equation}
is an isomorphism.

It remains to identify (5.2.4.6) with 5.2.4.2.  Let
$a,b\in H^{n-1}(X,\mathbb Q_\ell(j))$.  By Poincaré duality on $X$,
(5.2.4.6) is equivalent to nondegeneracy of the bilinear form
\begin{equation}\tag{5.2.4.7}
\begin{aligned}
H^{n-1}(X,\mathbb Q_\ell(j))\times
H^{n-1}(X,\mathbb Q_\ell(j))
&\longrightarrow H^{2n}(X,\mathbb Q_\ell(2j+1)),\\
(a,b)&\longmapsto a\,L_X(b).
\end{aligned}
\end{equation}
But by (5.2.4.3) and the projection formula,
\[
aL_X(b)=a\tau_*\tau^*(b)
=\tau_*\bigl(\tau^*(a)\tau^*(b)\bigr).
\]
This proves the equivalence of (5.2.4.7) and 5.2.4.2.

\subsection*{Corollary 5.2.5.}
\underline{Let $X$ be smooth and projective of dimension $n$, and let
$\tau:Y\hookrightarrow X$ be the inclusion of a smooth hyperplane section.
Suppose that $X$ satisfies (LV).  Then, for $q\geq n$, the map}
\[
\tau^*:H^q(X,\mathbb Q_\ell(j))
\longrightarrow H^q(Y,\mathbb Q_\ell(j))
\]
\underline{is surjective.}

\underline{Proof.}
By (5.2.4.3), for $i\geq0$ one has
\[
L_Y^{i+1}=\tau^*L_X^i\tau_*,
\]
and hence the commutative diagram
\begin{center}
\begin{tikzcd}[column sep=huge,row sep=huge]
H^{n-i}(X,\mathbb Q_\ell(j+1))
  \arrow[r,"L_X^i"]
  &H^{n+i}(X,\mathbb Q_\ell(j+i+1))\arrow[d,"\tau^*"]\\
H^{n-i-2}(Y,\mathbb Q_\ell(j))
  \arrow[u,"\tau_*"']\arrow[r,"L_Y^{i+1}"]
  &H^{n+i}(Y,\mathbb Q_\ell(j+i+1)).
\end{tikzcd}
\end{center}
This gives the required result.

\subsection*{Lemma 5.2.6.}
\underline{Let $f:F\to B$ be a smooth proper morphism with $B$ connected.
Let $\underline{\mathcal O}_F(1)$ be a relatively ample invertible sheaf on
$F$, and equip each geometric fiber $F_{\bar b}$ with the ample invertible
sheaf induced by $\underline{\mathcal O}_F(1)$.  If one geometric fiber
satisfies (LV), then so do all the others.}

\underline{Proof.}
Let
$\xi_F\in H^2(F,\mathbb Z_\ell(1))$ be the first Chern class of
$\underline{\mathcal O}_F(1)$.  This class and its powers
$\xi_F^i\in H^{2i}(F,\mathbb Z_\ell(i))$ induce, by cup product, the
homomorphisms
\begin{equation}\tag{5.2.6.1}
\xi_F^i:R^{n-i}f_*\mathbb Q_\ell(j)
\longrightarrow R^{n+i}f_*\mathbb Q_\ell(j+i),
\end{equation}
where $n$ is the relative dimension of the fibers.  By proper base change
[SGA 4, XII 5.2], all geometric fibers satisfy (LV) if and only if all the
maps (5.2.6.1) are isomorphisms.  By the specialization theorem
[SGA 4, XVI 2.2], the sheaves $R^if_*\mathbb Q_\ell(j)$ are twisted constant
sheaves.  Thus the maps (5.2.6.1) are isomorphisms if and only if they are so
over one given geometric point.

\subsection*{Corollary 5.2.7.}
\underline{Let $X$ be smooth and projective.  If one smooth hyperplane
section satisfies (LV), then every smooth hyperplane section does.}

\underline{Proof.}
The universal family of smooth hyperplane sections of $X$ satisfies the
hypotheses of Lemma 5.2.6.
